\documentclass[10pt]{article}
\usepackage[margin=20mm]{geometry}
\usepackage{amsmath,amssymb,mathtools,fontspec,unicode-math,xurl,hyperref}
\IfFontExistsTF{DejaVu Serif}{\setmainfont{DejaVu Serif}}{\setmainfont{Latin Modern Roman}}
\setmathfont{latinmodern-math.otf}
\setlength{\parindent}{0pt}\setlength{\parskip}{6pt}\sloppy
\title{Initial scientific review: M25-0043
Perverse filtrations and Fourier transforms}
\author{}
\date{25 October 2026}
\begin{document}
\maketitle
Abstract: The inspected Fourier-correspondence argument is consistent with its stated one-sided vanishing hypothesis. The supplied 2025 publisher version has a missing strictness argument in Corollary 5.3, used by the plane-curve application. This is a reproduced known issue: the authors’ subsequent arXiv v3 updates that proof. The strictness bridge in that repair is independently checked below. No additional substantive issue was established in the declared coverage; this review does not certify all 69 pages.\par
\section*{Identity, version and disposition}
The reviewed source is Davesh Maulik, Junliang Shen and Qizheng Yin, Perverse filtrations and Fourier transforms, Acta Mathematica 234 (2025), 1–69, DOI `10.4310/acta.2025.v234.n1.a1`. Printed and one-based PDF page numbers agree. The assigned file is `papers/0043/0043.pdf`, SHA-256\par
> e18c83ebeee24e45ba39a26307ab294d934
> 5667a4726deedff6ebfe8e1b230df.\par
The complete manifest and file bindings were verified at intake. Their parsing does not count as mathematical review. The later author manuscript is comparison evidence and has not replaced the publisher PDF.\par
Disposition. One known-issue reproduction, K01, concerns the proof of Corollary 5.3. Its dependency on a main application makes the gap Major-level for the frozen proof, but it is not reported as a new unresolved defect in the authors’ current manuscript. There are zero additional confirmed Major and zero additional confirmed Minor findings. No counterexample to the corollary or to Theorem 0.6 is asserted.\par
\section*{Scope, definitions and positive checks}
Text inspected: pp. 1–8, 12–16, 21–31, 40–42, 55–60 and 66–67 (35 of 69 pages). Actual source images inspected: pp. 1, 3, 7, 14, 15, 23, 24, 26, 27, 30, 31, 41, 42, 57, 58 and 60. The following are substantive checks, not a claim that every dependency on those pages was reconstructed.\par
\subsection*{Hypotheses and ranges}
The paper works over \(\mathbb C\) with rational coefficients. Its abelian fibrations have smooth irreducible total space and base, are proper and equidimensional, and admit the specified group action. Definition 1.2, pp. 15–16, includes inverse Fourier kernels, a convolution kernel with the requisite codimension bound, and full support of the perverse summands for the indicated fibration. These are hypotheses, not consequences of being any singular family of abelian varieties.\par
Theorem 0.1, p. 3, uses Fourier vanishing (FV). The compactified-Jacobian application in Theorem 0.2, p. 4, concerns integral projective planar curves, a section through the smooth locus, and a nonsingular total compactified Jacobian. Theorem 0.6, p. 7, asserts \(P\supset C\) for coprime \((r,\chi)\), \(r\geq3\), only in cohomological degrees at most \(2(r-2)\), not unrestricted \(P=C\).\par
With the paper’s unshifted direct-image convention, p. 14, \[P_kH^d(M)=\operatorname{im}\bigl[
 H^d(B,{}^p\tau_{\leq k+\dim B}R\pi_*\mathbb Q_M)\longrightarrow H^d(M,\mathbb Q)\bigr].\] The base-dimension shift is consequential in K01.\par
\subsection*{Riemann–Roch normalization and projectors}
The two Fourier classes on pp. 22–24 have deliberately asymmetric Todd factors. With \[T_{M^\vee\times_B M}
 =\operatorname{pr}_1^*T_{M^\vee}+\operatorname{pr}_2^*T_M-\pi^*T_B,\] the factors in the convolution cancel as required by Lemma 2.3. The proof uses the base-changed diagonal and flatness. Replacing its localized Chern character by an ordinary Chern character on an arbitrary singular product would be unjustified; the inspected proof does not do so.\par
Definition 2.5, p. 24, requires \(F_i^{-1}F_j=0\) for \(i+j<2g\) as relative Chow correspondences. No vanishing for \(i+j>2g\) is assumed. Set \(a_i=F_iF^{-1}_{2g-i}\). The proof on p. 26 has \(\sum_i a_i=1\) and \(a_i a_j=0\) for \(i>j\), since \(2g-i+j<2g\). For \(p_k=\sum_{j\leq k}a_j\), \[p_k=1p_k=\sum_i\sum_{j\leq k}a_i a_j
       =\sum_{i\leq k}\sum_{j\leq k}a_i a_j=p_k^2 .\] Thus \(q_{k+1}=1-p_k\) is its orthogonal complement. Only the stated one-sided vanishing is used. For \(g=0\) this reduces to \(p_0=1\).\par
The realization step on p. 27 uses full support and semisimplicity to exclude additional boundary-supported perverse summands. Generic agreement alone would not identify arbitrary derived morphisms; the split image and complementary summands permit the truncation identification here.\par
\subsection*{Multiplicativity and degree bookkeeping}
In the proof on pp. 28–30, a possibly nonzero component of the projected cup product must satisfy \[j_1\geq2g-k,\qquad j_2\geq2g-\ell,\qquad j_3\geq k+\ell+1.\] The convolution correspondence has degree at least \(-g\). Hence the total degree is at least \[(g-k)+(g-\ell)+(k+\ell+1-g)-g=1.\] The projected cup product is homogeneous of correspondence degree zero, so vanishes. This is a grading argument, without the missing opposite half of (FV). The Chern inclusion on p. 31 similarly uses \(2g-i+k<2g\) for \(i\geq k+1\).\par
For the subsequent motive splitting the preliminary projectors obey \(p_\ell p_k=p_k\) when \(\ell>k\). The replacement \(t_k=p_kp_{k+1}\cdots p_{2g}\) is nested on both sides: in \(t_\ell t_k\), successive \(p_i p_k\) collapse; in \(t_k t_\ell\), collapse the factors to the right down to \(p_\ell\). Thus \(t_\ell t_k=t_k t_\ell=t_k\) for \(k\leq\ell\), and \(t_k-t_{k-1}\) are orthogonal projectors. Commutativity of the original \(p_k\) is not assumed.\par
\subsection*{Supporting checks}
On pp. 40–42 the localized character of \(\psi^N\mathcal P\) scales in degree \(j\) by \(N^j\). Conditional on the support theorem used there, for fixed total degree below \(2g\) one obtains a finite relation \[\sum_{i+j=k}N^{j+\dim B}A_{ij}=0
 \qquad\text{for every positive integer }N.\] Taking as many distinct \(N\) as coefficients gives an invertible rational Vandermonde matrix after dividing the common nonzero power. Every \(A_{ij}\) therefore vanishes, even for an infinite-dimensional ambient Chow group. The Todd expansion has nonnegative index \(j'\) and replaces \(j\) by \(j-j'\leq j\), preserving the strict inequality. The preceding support theorem has not itself been reconstructed.\par
\section*{Problem 1: known missing strictness step (K01)}
\subsection*{Exact source and suspect inference}
Corollary 5.3, pp. 57–58, claims that restriction to the integral-curve locus \(W\) induces an isomorphism of perverse-filtered cohomology in degrees \(d\leq2r-4\), for every index \(k\). After triangle (5.5), p. 58 asserts low-degree ordinary and perverse bounds and claims an isomorphism of perverse cohomology sheaves for \(i\leq2r-4\).\par
The missing implication is strictness with respect to the whole perverse filtration in those ordinary cohomological degrees. Ordinary cohomological degree, absolute perverse-sheaf degree, and filtration index are different variables. The printed final inference does not supply that implication.\par
\subsection*{Complete check of the inference}
Let \(c=r-1\). Grant the geometric bound \(\operatorname{codim}_M Z\geq c\) used in the proof. The ordinary support bound and long exact sequence give an isomorphism \(H^d(M)\to H^d(M_W)\) for \(d\leq2c-2\): both adjacent support groups vanish. This proves the unfiltered assertion.\par
A filtered linear isomorphism need not be a filtered isomorphism: take \(V=W=\mathbb Q\), \(\phi=1\), with \(P_0V=0\), \(P_0W=\mathbb Q\) and \(P_1V=P_1W=\mathbb Q\), extended constantly above and by zero below. Then \(\phi(P_kV)\subseteq P_kW\) for every \(k\) and \(\phi\) is bijective, but \(\phi(P_0V)\ne P_0W\). This is a counterexample to the abstract inference, not a geometric counterexample to Corollary 5.3.\par
The perverse-sheaf range as printed also does not address every required truncation. Here \(N=\dim|rH|=r(r+3)/2\), and \(P_k\) uses cutoff \(k+N\). For \(r=3\), the printed bound is \(i\leq2\), whereas the \(P_0\) cutoff is \(9\). Even granting that low-degree perverse isomorphism does not yield the asserted filtered isomorphism.\par
An arbitrary proper direct image does not preserve all lower perverse bounds. For example, for \(f:\mathbb P^1\to\mathrm{pt}\), \(\mathbb Q_{\mathbb P^1}[1]\) is perverse but \[Rf_*(\mathbb Q_{\mathbb P^1}[1])=\mathbb Q[1]\oplus\mathbb Q[-1].\] This rules out one possible general justification; it does not show that the particular perverse bound in the paper is false. K01 rests on the missing strictness bridge.\par
The base-boundary count agrees with Lemma 5.2: if \(N_d=d(d+3)/2\), then \(N_r-N_a-N_{r-a}=a(r-a)\geq r-1\) for \(1\leq a<r\). Geometry transferring a base bound to the inverse image remains a dependency; that transfer is not proved by this polynomial identity.\par
\subsection*{Counterevidence, known revision and repair}
The arXiv record for `2308.13160v3`, dated 8 October 2025, records an update of the proof of Corollary 5.3. The revised proof combines the ordinary cohomology isomorphism with functoriality of a Deligne splitting, citing [Retained latex: \textbackslash{}cite[Lemma 3.3]\{dms\}]. That lemma and its proof were read on printed p. 937 of the primary source. The HTML rendering date is not treated as the manuscript version date.\par
Here is the strictness bridge independently completed. Choose an \(h\)-relative ample class \(\eta\); its restriction is relative ample for \(h_W\). Restriction is a graded ring map, preserves the perverse filtration by open-base-change naturality, and carries \(\eta\) to its restriction. The comparison lemma makes it preserve the first Deligne summands: \(\phi(G_kH^d(M))\subseteq G_kH^d(M_W)\). Fix \(d\leq2r-4\), where \(\phi\) is an ordinary isomorphism. For \(w\in G_kH^d(M_W)\), decompose its unique preimage as \(v=\sum_jv_j\) with \(v_j\in G_jH^d(M)\). The target direct sum implies \(\phi(v_j)=0\) for \(j\ne k\), so injectivity gives \(v_j=0\). Thus \(\phi\) is an isomorphism on every \(G_k\), and hence every \(P_k=\bigoplus_{j\leq k}G_j\).\par
This proves the missing bridge under the stated geometric and splitting inputs. It agrees with the later author proof and requires no all-degree ordinary isomorphism. It does not reprove all background geometry or the decomposition theorem.\par
\subsection*{Impact}
Theorem 0.6, p. 60, explicitly uses Corollary 5.3 to reduce to \(M_W\). Descending its perverse-degree bound requires the inverse restriction to preserve the filtration, exactly the missing point. The gap is core-linked even though repair is available. Frozen-proof severity is Major-level; reporting status is known-issue reproduction, later repaired, not a new current Major. Neither the abstract Fourier theorem nor the multiplicativity argument checked above is refuted.\par
\section*{Limits and conclusion}
No author code or empirical experiment was run. No separate computational repository or supplement was acquired in the manifest; absence is not a defect. Checks here are analytic algebra, degree bookkeeping and explicit examples.\par
Not performed: complete reconstruction of Arinkin’s autoduality and support estimates preceding Corollary 3.5; all localized intersection-theory foundations; twisted Jacobian and tautological constructions; every full-support and fiber-dimension result; the entire \(P=W\) reduction and its external dependencies. Section 5 was not certified in full.\par
Judgment | K01 is a known core-linked proof gap in the supplied version; no additional substantive issue established.\par
Repair | Strictness bridge verified under the stated geometric bounds and Deligne-splitting input.\par
Adoption | Authors’ v3 already updates the proof; no contact or new adoption request made.\par
Remaining | Uninspected arguments above and independent human scientific acceptance.\par
\section*{Supporting checks}
Existing Windows Poppler extracted the source text. Listed source pages were actually displayed; extraction and rendering alone are not review. No static author-code inspection is claimed.\par
No compiler executable was installed.\par
Human scientific adjudication remains pending.\par
9 D. Maulik, J. Shen and Q. Yin, Perverse filtrations and Fourier transforms, Acta Math. 234 (2025), 1–69. DOI: `10.4310/acta.2025.v234.n1.a1`. Assigned publisher PDF.\par
D. Maulik, J. Shen and Q. Yin, Perverse filtrations and Fourier transforms, arXiv:2308.13160v3, 8 October 2025. https://arxiv.org/abs/2308.13160v3. Versioned Corollary 5.3 and metadata consulted 25 October 2026.\par
M. A. de Cataldo, D. Maulik and J. Shen, Hitchin fibrations, abelian surfaces, and the \(P=W\) conjecture, J. Amer. Math. Soc. 35 (2022), 911–953. DOI: `10.1090/jams/989`. Lemma 3.3 and preceding description, pp. 936–937; primary author-hosted publisher offprint consulted 25 October 2026.\par
\end{document}
